\begin{table*}[!t]
\centering
\small
\renewcommand\tabcolsep{1.0pt} % column space
\begin{tabular}{cp{11.5cm}}
    \toprule
    \textbf{Locality Reasoning} & \textbf{Description} \\
    \midrule
    \multirow{3}{*}{\parbox{4cm}{\centering Comparison \\ \centering(\text{38.71\%})}}
    &  It focuses on evaluating and contrasting information presented in tables, figures, or other data formats. To answer questions of this type, one must analyze and compare specific subjects or variables within the given data set. \emph{Example: Which model shown in the figure consistently better than MPT-7B-Instruct but consistently worse than LLaMA-30B}. \\
    \midrule
    \multirow{3}{*}{\parbox{4cm}{\centering Data Extraction \\ \centering(\text{33.47\%})}}
    &  It directly retrieves specific information from table or figure. This approach focuses on pinpointing exact data points or details. \emph{Example: Which model demonstrates score of 21.073 in 10-shot prompting}\\
    \midrule
    \multirow{3}{*}{\parbox{4cm}{\centering Location  \\ \centering(\text{15.29\%})}}
    &  It centers on pinpointing spatial or positional information from table or figure. This involves identifying either relative or absolute locations, like the placement of items in a figure or row information in a table. \emph{Example: Which dataset lies on the second left in the figure?}\\
    \midrule
    \multirow{3}{*}{\parbox{4cm}{\centering Visual Understanding \\ \centering(\text{12.53\%})}}
    & It emphasizes on the visual information from the figure such as colors, shapes, and marker types. This approach involves closely analyzing these visual cues to extract relevant information. \emph{Example: What is the model represented with a blue line in the figure?}\\
    \toprule
    \textbf{Detail Reasoning} & \textbf{Description} \\
    \midrule
    \multirow{3}{*}{\parbox{4cm}{\centering Conceptual Understanding \\ \centering(\text{21.62\%})}}
    &  Evaluate knowledge of essential concepts, basic theories, and critical definitions related to the subject. \emph{Example: What does the hypernetwork in the proposed Hyperdecoders approach generate?}\\
    \midrule
    \multirow{3}{*}{\parbox{4cm}{\centering Methodological Analysis \\ \centering(\text{21.01\%})}}
    & Examine and assess the research methodologies and experimental frameworks employed in studies, with an emphasis on their efficacy and constraints. \emph{Example: What potential application of the Hyperdecoder approach is suggested by its performance on long-context out-of-domain datasets in the MRQA evaluation?} \\
    \midrule
    \multirow{3}{*}{\parbox{4cm}{\centering Results and Data Interpretation \\ \centering(\text{19.49\%})}}
    & Analyze statistical data, graphs, and tables, focusing on deriving significant insights and conclusions from quantitative and visual information. \emph{Example: In the experimental results for the GLUE benchmark using $T5_{large}$ v1.1 + LM as the underlying model, which model configuration achieved the highest average score across tasks?} \\
    \midrule
    \multirow{3}{*}{\parbox{4cm}{\centering Implications and Inferences \\ \centering(\text{19.21\%})}}
    & Infer wider implications and practical uses of study outcomes, concentrating on the extensive impact and prospective significance of the results. \emph{Example: What is the primary purpose of introducing domain prompts in ELLE?}\\
    \midrule
    \multirow{3}{*}{\parbox{4cm}{\centering Critical Analysis \\ \centering(\text{18.66\%})}}
    & Assess the study’s reasoning, robustness of evidence, and validity of conclusions critically, with a focus on logical consistency and the support of empirical data. \emph{Example: What is the primary goal of ELLE as described in the paper?} \\
    \bottomrule
\end{tabular}
\caption{Definitions and proportions of reasoning categories in \data.}
\label{tab:reasoning definition}
\end{table*}